% Original topic-level problem, solution and MCQ reasoning.
\subsection{বাস্তব সমস্যা: পুরো কলেজে এক ধরনের মাধ্যম নয়}
\chTwoPair{\textbf{সমস্যা:} একটি কক্ষে ১৫টি কম্পিউটার আছে। দুটি ভবনের দূরত্ব ৭০০ m। শিক্ষার্থীরা ল্যাপটপ নিয়ে স্থান বদলায়। এসব কাজে কম খরচে নির্ভরযোগ্য সংযোগ দরকার।\par
\textbf{সমাধান:} কক্ষের কম্পিউটারগুলো উপযুক্ত Ethernet UTP ও সুইচ দিয়ে যুক্ত করা যায়। ভবনের প্রধান সংযোগে ফাইবার এবং চলমান ল্যাপটপে Wi-Fi/AP ব্যবহার করা যায়। সাধারণ ১০০ m তামার চ্যানেল দিয়ে একটানা ৭০০ m সংযোগ পরিকল্পনা করা ঠিক নয়।}{\textbf{ব্যবহার ও যুক্তি:} স্থির যন্ত্রে কেবল সংযোগ সহজে নিয়ন্ত্রণ করা যায়। দূরের প্রধান সংযোগে ফাইবার বেশি বহনক্ষমতা ও EMI থেকে সুরক্ষার সুবিধা দেয়। Wi-Fi দিয়ে চলাচলের সময়ও যোগাযোগ করা যায়। সব অংশের ডেটা হার সমান নাও হতে পারে।\par
\textbf{মাধ্যম নির্বাচনের শর্ত:} দূরত্ব, বহনক্ষমতা, বাধা বা EMI, চলাচল, বাজেট ও রক্ষণাবেক্ষণ একসঙ্গে বিবেচনা করতে হয়। শুধু দ্রুত বা সস্তা দেখে মাধ্যম বেছে নেওয়া যায় না।\par
\textbf{একই Network-এ ভিন্ন Medium:} একই নেটওয়ার্কের পৃথক সংযোগে তারযুক্ত ও বেতার মাধ্যম ব্যবহার করা যায়। টপোলজি ও মাধ্যমের শ্রেণিবিভাগ আলাদা।\par}

দিকনির্দেশিত মাইক্রোওয়েভের সরাসরি দৃষ্টিপথ বা LOS এবং পুনঃপ্রেরণের ব্যাখ্যা বেতার মাধ্যমের অংশে আছে। দূরের সংযোগ পরিকল্পনায় এসব শর্তও মানতে হয়।
